\documentclass[11pt,a4paper]{ltjsarticle}
\usepackage{luatexja-fontspec}
\IfFontExistsTF{Noto Serif CJK JP}
 {\setmainjfont{Noto Serif CJK JP}\setsansjfont{Noto Sans CJK JP}}
 {\IfFontExistsTF{Harano Aji Mincho}
   {\setmainjfont{Harano Aji Mincho}\setsansjfont{Harano Aji Gothic}}
   {\setmainjfont{Hiragino Mincho ProN}\setsansjfont{Hiragino Kaku Gothic ProN}}}
\setmonofont{Latin Modern Mono}[Scale=0.88]
\usepackage{amsmath,amssymb,amsthm,mathtools,mathrsfs}
\usepackage[a4paper,margin=24mm,headheight=16pt]{geometry}
\usepackage{booktabs,longtable,array,enumitem,needspace,xcolor}
\usepackage[unicode,bookmarksnumbered,linktoc=all]{hyperref}
\hypersetup{colorlinks=true,linkcolor=blue!45!black,citecolor=green!30!black,urlcolor=blue!50!black}
\renewcommand{\contentsname}{目次}
\renewcommand{\refname}{参考文献}
\renewcommand{\proofname}{証明}
\newtheoremstyle{japanese}{8pt}{6pt}{\normalfont}{}{\bfseries}{．}{0.4em}{}
\theoremstyle{japanese}
\newtheorem{theorem}{定理}[section]
\newtheorem{proposition}[theorem]{命題}
\newtheorem{lemma}[theorem]{補題}
\newtheorem{corollary}[theorem]{系}
\newtheorem{definition}[theorem]{定義}
\newtheorem{remark}[theorem]{注意}
\numberwithin{equation}{section}
\setlist{leftmargin=2em,itemsep=3pt,topsep=4pt}
\setlength{\parindent}{1em}
\setlength{\parskip}{0.2em}
\setlength{\emergencystretch}{2em}
\allowdisplaybreaks[1]
\newcommand{\rk}{\operatorname{rank}}
\newcommand{\HH}{\operatorname{HH}}
\newcommand{\versiondate}{2026年11月19日・第4版}
\setcounter{tocdepth}{1}

\title{研究ノート：非AS性・Clifford orderと超曲面障害}
\author{}
\date{2026年11月24日／non-as-v11}
\begin{document}
\maketitle
本書は現在有効な補助知識を集める。最初に研究方針 \path{00_research_plan.md} を読む。
完成した全証明は一本の公表用論文 \path{papers/nc-cubic-non-as.tex} に収める。
本ノートは証明を重複させず、主論文の各節を参照する。

\section{最優先課題：OU三次型と標準超曲面型の逆包含}
\label{sec:ou-converse-notes}
主論文の定理\texttt{thm:ou-weak-cubic}・系\texttt{cor:ou-weak-unpointed}に、
OUの箙型による三次del Pezzo型の逆包含を否定する証明がある。
内在的候補$\mathfrak I_3$との共通部分での逆包含はなお未解決。
このより狭い問いの決着を最優先とし、既存定理の独立監査も継続する。

\subsection*{弱del Pezzo境界の機構}
平面楕円曲線上の六点のうち$p_1,p_2,p_4$のみが共線。
AOU (4.1) の第一ブロックでは$p_1,p_2,p_3$も$p_4,p_5,p_6$も
共線でなく、全六点がconicにも載らないことが重要。
従って共線性はvery strong helixを破壊しない。
反標準収縮上のheartは$\operatorname{coh}(Y,f_*\mathcal End T)$であり、
通常の$\operatorname{coh}X$と混同しない。
$qgr$の同定、箙の型（逆Serre bimoduleによる直接計算）、
Serre移動、AZ2からの局所射影性は主論文を参照。
$(-2)$曲線の収縮によって通常の点対象がheart内で
二つの零次元対象への拡張となる。この分解が任意のample対象に
導来rankの下界を与える。これはHochschild障害ではない。
$K_X^2=3$であっても、heartの反標準Hilbert多項式の次数規格化は
$3r^2$となる。従って導来圏の数値degreeだけで内在的degreeを判定できない。

一般化の候補：$D^b$同値により滑らかな曲面上の点対象がheartに入り、
反標準移動で固定される非零対象への長さ$m$のfiltrationを持つ場合、
AZ-ample対象のvirtual rankが$m$以上となる。
中央の豊富束による局所射影性（あるいは点支持対象に対する
高次Ext消滅）が必要であり、任意の非可換heartに無条件で適用しない。
複数の収縮曲線、ADE型、異なるfoundationについてこの条件と
最小rankの実現を調べる価値がある。

\subsection*{肯定的な範囲と適用境界}
主論文の定理\texttt{thm:recognition-weak-commutative}により、
可換次数三弱del Pezzo導来型は全て判定済み：
内在的候補なら滑らかな可換三次曲面である。
個別のfoundationや導来rank障害に依存しない。
中心上有限の反標準環の一般点での次元一とtorsion-free性が核心であり、
中心上有限性がない場合へ無条件に拡張してはならない。

主論文の命題\texttt{prop:helix-heart-upgrade}と
系\texttt{cor:aou-abelian-upgrade}により、
AOU Theorem 9.4 の比較が適用される範囲では
自由対象の対応を起点として標準heartを認識できる。
導来同値だけから任意の標識の同一性を結論しない。
VdBの爆発定理が適用されることと、AOUのhelix比較が
適用されることも区別する。

\subsection*{既約種数一曲線型の直接認識と境界の計算}
主論文の定理\texttt{thm:ou-integral-elliptic-recognition}は
滑らかな楕円曲線・節点・尖点三次曲線上の次数$0,1,2$各三本の線束と
標準noetherian qgr/導来helix比較を持つOU型が内在的三次候補でもあるとき、
与えた標識の反標準環を標準AS4三次商と認識する。
爆発の一般位置を仮定しない。
特異曲線の場合に$\operatorname{Perf}E$と$D^b(\coh E)$を区別する。
OU spherical functorの三角は前者を用い、右随伴$J$は九線束からの
導来Homにより後者にも延長する（$\Lambda$の有限大域次元を使用）。
Burban--Kreu{\ss}ler Cor.~5.8 と二段の評価核の計算により、
九本のtwistの積の逆を$[2]$した$h$は
次数三の線束による捩りと曲線自己同型であり、偶数シフトは零。
行列式を取るべきGrothendieck群は$K_0(\operatorname{Perf}E)$であり、
特異曲線の$G_0(E)$で代用しない。
この行列式から$\alpha^*$が全次数零の線束に恒等であること、
ATVのregular triple条件$(\alpha^*-1)^2[\mathcal L]=0$が従う。
$J$の完全性と像のGK次元一以下を確認し、反標準三角と
$\mathcal O$のtorsion-free性、AZ1/AZ2で制限線束を認識する。
正則中心一次元の商は次数3 THCR、二次包絡は中心拡大判定による。
指定の曲線データだけで内在的条件や全OU型の被覆を結論しない。

\subsection*{AZ2からの局所自由性と可約曲線への帰着}
主論文の補題\texttt{lem:ou-az2-singular-free}は連結被約射影曲線全般で
幾何的$h$に対するAZ2から全特異点での局所自由性を導く。
全射$\mathcal O_E\to k(p)$と$p$の周期を使い、
剰余体値の全線形汎関数が局所環値の写像に持ち上がることから、
Nakayamaの補題を二度適用する。被約性と有限特異点集合が必要。
次に$M$はベクトル束と滑らかな点に台を持つ有限長層に分解する。
命題\texttt{prop:ou-reduced-curve-rank}では、導来随伴$F\dashv J$、
反標準三角、$Fs\simeq hF$、完全conservativeな$J$と像GK次元一以下、
$\deg N=3$を仮定している。
AZ2を曲線側へ移し、AZ1とHilbert差$3n$から
$3n\ge 3nr^2+r\ell$で$r=1,\ell=0$となる。
これは被約可約算術種数一曲線にも使える。

命題\texttt{prop:ou-reduced-cubic-geometric}ではさらに$h$が
九線束のtwist積に由来すると仮定し、全次数零線束への作用から
成分置換を排除する。AZ1が曲線側に移るので
$\deg N|_{E_i}>0$（各成分）、全次数三を得る。
被約平面三次曲線では真の部分曲線とその補曲線の交叉長は二。
この値と正多重次数から長さ二の部分スキームを分離でき、
very ample性とregular triple性が従う。
従って可約被約曲線でも上記仮定の下で標準AS4三次商。
$R/zR$が整域であることはこの中心拡大判定に不要。
幾何的でampleな$h$と標準helix比較があれば曲線heart認識は
既約の場合と同じ方法で示せる。次の多重次数判定によりこの範囲を特定した。
非被約曲線では特異点集合が有限でないため、上の局所補題を流用しない。
Burban--Kreu{\ss}ler Cor.~5.8の可約曲線への無条件拡張も行わない。

\subsection*{可約被約三次曲線の多重次数判定}
主論文の定理\texttt{thm:ou-reduced-md-criterion}と
系\texttt{cor:ou-reduced-md-recognition}を参照。
九線束のブロック和を$A,B,C$とし、
$\delta_{\rm md}=A+C-2B,\ w=5B-3A-2C$とおく。
この$\delta_{\rm md}$は楕円標識の$\Pic^0$の元$\delta$とは異なる。
$K_0(\Perf E)$のdeterminantを用い、仮想階数$r$、全次数$d$、
多重次数$x$に対する反復作用は
\[
 h^n:(r,d,x)\longmapsto
 \left(r,d+3nr,x+nrw+
       \left(nd+\frac{3r n(n+1)}2\right)\delta_{\rm md}\right).
\]
$\delta_{\rm md}=0$かつ$w_i>0$が幾何的ample移動の必要十分条件。
十分性は多重次数零線束のtwist共役を点層・部分正規化・FM核で認識し、
同じ曲線上の可換三次曲面由来の参照配置と比較する。
Opper (JEMS 28 (2026), Prop.~8.40, Cor.~8.41) は節点cycleの比較に関係するが、
接触・共点も含む今回の証明はその分類に依存しない。
AOU (4.1)とOU (4),(5) を参照配置に適用する。
必要条件$-1\leq y_i-x_i\leq d+1$から有限列挙を行い、
成分を区別しブロック内をソート、共通捩りを正規化した規約で、
2成分180型中6型、3成分15012型中9型が上の条件を満たす。
すべて巡回変更で古典参照型に帰着。
\path{computations/verify_ou_reduced_multidegrees.py} を実行済み。
これは実際のOU圏や線束配置の同型類の個数ではない。
新証明の全独立監査・組版・新規性確認は未実施。

\subsection*{標識変更の独立な結果}
命題\texttt{prop:ou-marking-walls}を参照。
第一ブロック内の六順序でVdBの移動後一般位置を判定すると、
99 primitive行（18行は1順序、81行は2順序）を得る。
その3975の有理二次元行空間のうち全順序を覆う30空間は
すべて$\delta=L_1-L_0$を含む。
従って$\delta$非torsionなら全順序失敗は余次元3以上。
再現コード \path{computations/verify_ou_decuple_markings.py} の
実行済み記録を保持する。今回の全計算再実行とは区別する。
単一の混合三点等式の一般点は別の順序で回復できる。
AOU Thm.~9.4の全仮定をVdB一般位置のみと同一視しない。
非torsionの場合の三点一般位置条件と関係係数非消失条件は異なる
（主論文の命題\texttt{prop:decuple-zero-coefficients}と
注意\texttt{rem:decuple-distinct-boundaries}）。

\subsection*{内在的条件による可約曲線の残余型の排除}
命題\texttt{prop:ou-residual-delta-obstruction}および
定理\texttt{thm:ou-residual-geometric-ample}を参照。
標準noetherian qgr/導来helix比較を仮定する。
反標準自然変換$\epsilon:s^{-1}\to\mathrm{id}$がheart上のどこかで
非零なら、AZ1とrank-one torsion-free性から全ての$s^a\mathcal O$上で
非零かつ単射となる。これは$h$の幾何性に依存しない。
$\epsilon_{JM}=0$は$F\epsilon=0$と随伴から従うので、
$s$の$t$-exact性により$J(M)$の各cohomologyも$\epsilon$で消される。
そのような対象の切断数は、AZ1による有限直和からの全射と
AZ2によるHom全射、$\dim R_n-\dim R_{n-1}=3n$から$O(n)$。
この上界はhelixの切断数にも移る。$M=\mathcal O_{E_a}$について
\[
 \chi_E(h^{-n}L_i,\mathcal O_{E_a})
 =1-(\operatorname{mdeg}h^{-n}L_i)_a
 =-\tfrac32n^2(\delta_{\rm md})_a+O(n).
\]
Euler数の線形成長と比較して$\delta_{\rm md}=0$を得る。
$J$の$t$-exact性も曲線側の幾何性もこの議論には不要。
なお反標準三角が存在することから
$\epsilon$がheart上のどこかで非零であることは自動ではない。

$h$が幾何的ならば、very strong性から$w_a\geq0$、
$\epsilon_{sP_i}\ne0$を得る。全閉点の$Jk(p)$はheartに属し、
十分負のhelixからの切断数は9である。
三ブロックEuler行列のSerre作用は$(S-1)^3=0$であり、
十分負のhelixのExt消滅・生成性により
切断数が有界なheart対象は有限長となる。
$M=F\mathcal O$は完全複体で、点層へのHomとSerre双対性から
Tor振幅$[-1,0]$。$H^0(M)$の特異点での自由性は
$J$がheartを全体として保つと仮定せず、
十分正の多重次数の$A_p$についてのみ
$J(A_p)\twoheadrightarrow Jk(p)$を用いてAZ2から導く。
$H^0(M)$と$M$の全成分で一定のrank、$w$の正の成分上での
$H^{-1}(M)$の消滅判定により$H^{-1}(M)=0$。
Hilbert差とtorsion長による不等式から$M$は線束となる。
その結果、$Jk(p)$は単純であり、滑らかな点について
$H^0(FJk(p))=k(p)$なので相異なる点が相異なる単純対象を与える。
$w_a=0$なら成分$E_a\simeq\mathbb P^1$上の十分正の固定次数の線束の
押出し$X_m$について$JX_m$は有限長であり、全滑らかな$p\in E_a$の
$Jk(p)$がその単純商となる矛盾を得る。
従って幾何的$h$は内在条件の下で自動的にample。
先の標識付き超曲面認識の適用にはこれで十分である。

多重次数の有限候補数を$\delta_{\rm md}=0$に限定すると
2成分16型（$\min w>0,=0,<0$に分けて6,6,4）、
3成分168型（9,81,78）である。
\path{computations/verify_ou_reduced_multidegrees.py}を拡張して再実行した。
この個数は実現した線束配置やOU圏の同型類の数ではない。
正の6/9型は既認識、残りの幾何的$h$は$I_3$から排除されたが、
非幾何的な$h$の排除とは異なる。

\subsection*{残るOU型への具体的入口}
最優先は曲線表示の反標準自然変換$\epsilon$のheart上の非零性を
一般に証明するか、全てのheart対象で零となる場合を排除することである。
具体的には$\epsilon_{P_i}$をdirected algebra上の
$J(L_i)$の二つのコホモロジーの拡大類として計算する。
$\Lambda$の射影$P_i$と対応する入射$I_i$を用いると、
反標準三角は$I_i[-2]\to P_i\to J(L_i)$となり、
調べる写像は$\operatorname{Ext}^2_\Lambda(I_i,P_i)$の元である。
一つでも非零なら$\delta_{\rm md}=0$を導ける。
この拡大類の非零性の一般証明はまだない。

次に$\delta_{\rm md}=0$でなお非幾何的なtwist積の振幅と
成分別の作用を解析し、幾何性・別表示・内在条件による排除を調べる。
AOUの$M_{\mathrm{rel}},M_\chi$における曲線表示の像・補集合と
OU三次型の全定義範囲を区別し、非被約・非曲線型を調べる。
非被約の場合に特異点有限性のAZ2補題を流用しない。
標準noetherian qgr/導来比較は定理の仮定である。
OU arXiv:2007.07620v1 pp.~1--2, (3)--(5) と
OU arXiv:2508.04986v1 Def.~1.1, \S3, Thm.~6.4を今回確認した。
後者の曲線データとの比較はbirationalであり、全OU型の被覆ではない。
独立監査・網羅的新規性確認・組版は残る。

\section{候補定義の階層と商スタックの検査}
主論文の内在的候補の定義
に対応する内在的候補$\mathfrak I_d$に対し、次数三では
\[
 \mathfrak H^{\mathrm{hyp,OU}}_3
 \subsetneq\mathfrak H^{\mathrm{hyp}}_3
 \subsetneq\mathfrak I_3
\]
が示される。ただし最初のクラスは超曲面型とOU型の\emph{共通部分}であり、
OU型全体については弱del Pezzoの反標準rank障害により逆包含が偽となる。
ただし内在的候補にも属するOU型での逆包含は未解決である。
第一の真性はrank一の符号型族の$\HH_{\pm1}=2$、第二はClifford例の
$\dim Q(R)_3=19$とunpointed非表示による。
量子Fermatの同じHochschild障害も独立の結果だが、包含の証明では符号型を採用している。比較の証明は主論文の冒頭と各定理を参照する。

$\mathfrak X=[\mathbb P^2/\mu_3]$（対角重み$(1,-1,0)$）では
$s=-\otimes\omega_{\mathfrak X}^{-1}$であり、
$H^0(\mathbb P^2,\mathcal O(3n))^{\mu_3}$から
$\mathbb C[U,V,W,T]/(UVW-T^3)$が得られる。
固定点のstabilizerは$\mu_3$、接空間の非自明な二指標は
$\chi,\chi^{-1}$であるため、粗い反標準模型は三つの$A_2$特異点を持つ。
構造層の反標準twistは固定点で自明指標しか持たずAZ-ampleでない。
一方、正則表現とのテンソル$E$はrank三、AZ-ample、
$\operatorname{End}(E)=\mathbb C^3$、
$\dim\operatorname{Hom}(E,s^nE)=3\binom{3n+2}{2}$。
本例はrank条件を緩める際にはconnectedness・exceptionality・
反標準模型の特異性を個別に検討すべきことを示す。
詳細な説明・証明は主論文の注意\texttt{rem:mu3-stack}を参照する。

\subsection*{二次包絡の必要十分認識}
主論文の命題\texttt{prop:quadratic-envelope-recognition}は、
一次生成と標準三次Hilbert級数の下で$Q(R)$が14641型noetherian
AS四次元正則環であることと、$R=Q(R)/(g)$（$g$は三次の正則正規元）
であることの同値を示す。証明は次数三で核が1、次数四で核が4、
従って$gB_1=B_1g$という有限次数の比較による。
整域性はLevasseurの定理（文献\texttt{Levasseur1992}, Theorem 6.2(ii)）に依存する。
これは分類稿の認識定理を本稿中で証明し直したもので、
広い内在的候補で包絡正則性が自動だと主張しない。

\subsection*{商スタックのrank障害の精密化}
反標準束が慣性群$\mu_3$の三つの固定点ファイバーで自明指標を持つため、
AZ-ampleなtorsion-free対象$F$の二重双対$F^{\vee\vee}$は
各固定点で三つの既約指標を全て含まなければならず、rankは少なくとも3。
証明に必要なのは、滑らかな二次元スタックで
$Q=F^{\vee\vee}/F$が零次元、
線束$M$に対して
$\operatorname{Ext}^1(Q,M)=0$なので
$\operatorname{Hom}(F^{\vee\vee},M)\cong\operatorname{Hom}(F,M)$
であることと、AZ1の全射を各慣性群の表現で調べること。
正則表現$E$（rank3）がこの境界を達成する。
主論文の命題\texttt{prop:mu3-minimal-rank}に完全な証明がある。

\section{三例の共通構造と次数一・二の補助知識}
三例とも $\mathcal A=\operatorname{Cl}_0(\mathcal V,q,L)$、
$\mathcal P=\operatorname{Cl}_1(\mathcal V,q,L)$ とし、右加群を用いる。
反標準自己同値は $T=-\otimes_{\mathcal A}\mathcal P$ である。
$a\in\operatorname{Hom}(F,T^mF)$、$b\in\operatorname{Hom}(F,T^nF)$ の orbit 積は
$ab=T^n(a)\circ b$ とする。この規約で Clifford の通常の積と一致する。
四対象の合成不変量の次元には、全ての矢を同時に逆順で記す選択は影響しないが、
証明書と本文の文字順序は統一する。

\begin{center}
\begin{tabular}{c|c|c|c|c}
\toprule
$e$ & $Z$ & $L$ & $D$ & $g(D)$\\\midrule
1 & $\mathbb P^2$ & $\mathcal O(1)$ & 五次曲線 & 6\\
2 & $\mathbb P^1\times\mathbb P^1$ & $\mathcal O(1,1)$ & $(3,3)$ 曲線 & 4\\
3 & $\mathbb F_1$ & $C+2f$ & $3C+4f$ & 3\\\bottomrule
\end{tabular}
\end{center}
三例の全てで $\det\mathcal V=\omega_ZL^2$、
$\mathcal P^{\otimes_{\mathcal A}2}=\mathcal A L$、
$\omega_{\mathcal A}=\mathcal P^{-1}$ となる。
smooth proper 性、Serre 双対性、rank 一の一般点、AZ ampleness と
負次数を含む切断の条件は、それぞれの論文の構成節で証明している。

\subsection{標準重みと別の ambient を区別する}
次数二の反標準環は最小生成次数 $(1,1,1)$ の AS 正則四次元代数の
中心四次超曲面であり、その Hilbert 級数の分母は
$(1-t)^3(1-t^2)$ である。
分母の $1-t^2$ は、独立な次数二生成元の存在を意味しない。
次数一では最小生成次数 $(1,1,3)$ の AS 正則四次元代数と中心六次式を用い、
分母は $(1-t)^2(1-t^2)(1-t^3)$ となる。
これらの PBW・Ore 表示と正則性の全証明は、主論文の次数一・二部分の
「反標準環の表示と AS-regular ambient」にある。
指定した重みの非表示を、全ての AS 正則 ambient に対する非表示へ強めてはならない。

両例とも $R_1\otimes R_1\to R_2$ が全射である。
これに対し標準重みの ambient は独立な次数二生成元を一つ持ち、
四次・六次式で割ってもその生成元は消えない。
任意のアーベル圏同値で自由対象の像を分類した後、
この三対象の合成写像の全射性を比較する。
次数三では $\dim R_2=10$ が標準 cubic と一致するため、
四対象の二次関係の三次帰結まで比較する必要がある。

\subsection{次数一の接触 conic}
次数一論文の判別式
\[
 \Delta=-u^5+2u^2vw^2+2uv^3w-2v^5-w^5
\]
に $(u,v,w)=(s^2,st,t^2)$ を代入すると
\[
 \Delta(s^2,st,t^2)=-(s^5-t^5)^2
\]
となる。従って $uw=v^2$ は判別五次曲線の五点接触 conic である。
$s^5=t^5$ の五点は標数零で相異なり、各交点の重複度は二である。
この恒等式は \path{scripts/nc_dp1_quintic_checks.py} で確認する。
既知の $F_{5-}^{2}$ family の theta 特性との比較に使えるが、
非表示定理の証明には必要ないため本ノートに保持する。

\subsection{次数一・二の対象分類と Hochschild homology}
局所射影 division rank 一の exceptional 対象は、次数二で
$T^i\mathcal A\otimes\mathcal O(j,k)$、次数一で $T^n\mathcal A$ に尽くされる。
Euler 行列の行列式はそれぞれ $16,1$ である。
分岐一般点での局所自由性、reduced norm による行列式の平方条件、
数値類から実際の同型へ移る議論を省略しない。
一般点の中心上の rank だけでは分岐点での二つの射影型の重複度は決まらない。

$HH_{-1},HH_0,HH_1$ の次元は次数二で $(3,4,3)$、次数一で $(5,3,5)$。
非零次数は判別二重被覆の Prym の接・余接空間として記述できる。
conic bundle の Chern class と Euler 数による独立計算も主論文の次数一・二部分に収める。
これは射影圏の標準 dg enhancement の HH であり、affine graded 環 $R$ の HH ではない。
これらの例自体が AS 超曲面表示と非零 $HH_{\pm1}$ の両立を示している。

\section{次数三の構成の選び方と数値}
quaternion order $\mathcal A$と可逆奇Clifford bimodule $\mathcal P$について
$\mathcal P^2=\mathcal A\otimes L$、$\omega_{\mathcal A}=\mathcal P^{-1}$
を実現するには、三変数二次束について
$\det\mathcal V=\omega_ZL^2$と置くのが便利である。
このとき$T^2=-\otimes L$であり、偶数次数の切断数の先頭項は
$\rk_Z(\mathcal A)L^2m^2/2=2L^2m^2$となる。
内在的次数を三にするには$L^2=3$が必要である。

$Z=\mathbb F_1$で$L=C+2f$を取ると$L^2=3$であり、
$\omega_ZL^2=\mathcal O(f)$である。
$\mathcal V=\mathcal O^2\oplus\mathcal O(f)$はこの条件を最も単純に実現する。
判別因子は
\[
D=3L-2c_1(\mathcal V)=3C+4f=-2K_Z-L.
\]
特に$K_Z+D=C+f$で種数は三、orderの標準因子は$-L/2$となる。
$K_{\mathcal A}^2=3/4$と反標準Hilbert次数$3$を混同しない。

固定する形式は
\[
q=\begin{pmatrix}cs^2+dt&ct^2&d\\ct^2&ct^2+ds&0\\d&0&c\end{pmatrix},
\quad
\deg(s,t,c,d)=(f,f,C,C+f).
\]
以下の到達点は全てこの明示式についてのものであり、
任意の形式や一般のパラメータに関する無条件な主張に読み替えない。

\section{主張の状態と証明の所在}
\begin{longtable}{p{.36\linewidth}p{.55\linewidth}}
\toprule
主張&状態・証明\\\midrule
判別曲線が滑らか・連結・種数三&
主論文の次数三部分の判別曲線の補題。四チャートの特異点イデアルは$\mathbb Q$上で$[1]$。
コードを実行済み。\\
一般点のquaternion algebraがdivision&
主論文の次数三部分の Clifford order の構成。
判別二重被覆の2-torsion線束$\eta$の解像度から$H^0(\eta)=0$を証明。\\
内在的三次曲面の全公理&
論文の構成節。局所order、Serre関手、ample pair、例外性、
Hilbert関数、smooth proper性の証明を記載し点検。\\
$R$の次数一生成・左右noetherian・domain&
主論文の次数三部分の反標準環の生成の命題。Cox単項式と一般点のdivision性による証明。\\
$\dim Q(R)_3=19$&
主論文の次数三部分の二次包絡の命題。
六関係から得る整数行列について、四十五次小行列式$512$と
独立な三本の整数核ベクトルを計算・検証。\\
局所射影rank一exceptional対象の分類&
論文の同名節。数値的群rank四、行列式の平方条件、Euler整数性、
数値類からの同型への昇格。\\
構造対象を忘れたgraded AS4三次商との非同値&
論文の最終定理。自由対象の像の局所射影性・rank一性を導き、
四対象の合成図式の十九次元と二十次元を比較。\\
任意の非周期的$\mathbb Z$-algebraの排除&
未完成。次節の追加仮定の下では同じ合成不変量で排除できる。\\
新規性&
未確認。ramification型の存在は既知。非表示定理の初出を今後照合する。\\
\bottomrule
\end{longtable}
幾何学・圏論の証明の点検と、計算コードの実行は別の検証である。
定理証明支援系による形式化は実施していない。

\section{障害の内在化と一般$\mathbb Z$-algebraの境界}
対象$F$について$E_i=T^iF$、$V_i=\operatorname{Hom}(E_i,E_{i+1})$、
$W_i=\ker(V_i\otimes V_{i+1}\to\operatorname{Hom}(E_i,E_{i+2}))$と置く。
\[
D_3(F)=
\frac{V_0\otimes V_1\otimes V_2}
{W_0\otimes V_2+V_0\otimes W_1}
\]
の次元は、四対象の合成図式だけから定まり、対象の同型の選択に依存しない。
本例の分類された全ての候補$F$で$\dim D_3(F)=19$となる。

connected gradedなAS三次商$B$では、Nakayama自己同型があり得ても、
$S(\pi B(i))\simeq\pi B(i-1)[2]$が対象として成り立つ。
四対象を個別に自由対象列と同定すれば、合成図式は通常の積の図式になり、
$\dim D_3=20$を得る。
$T=(1)$という関手の等式や、反標準orbit環と$B$の無条件な次数付き同型は不要である。

\begin{proposition}[追加仮定の下での拡張]\label{prop:conditional-Z}
$\mathbb Z$-algebraの三次超曲面圏$\mathcal B$がSerre関手$S_{\mathcal B}$を持ち、
$T_{\mathcal B}=S_{\mathcal B}^{-1}[2]$が完全自己同値であるとする。
さらにtautological対象列$(\mathcal B_i)_{i\in\mathbb Z}$があり、
以下を仮定する。
\begin{enumerate}
\item $\mathcal B_0$はexceptionalで、$T_{\mathcal B}\mathcal B_i\simeq\mathcal B_{i+1}$。
\item 自由対象間のHomと合成は次数付き成分に飽和し、
$\dim\operatorname{Hom}(\mathcal B_0,\mathcal B_n)=1+\frac32n(n+1)$ ($n\ge0$)。
\item 任意のcoherent対象$M$について
$\operatorname{Ext}^1(\mathcal B_{-n},M)=0$ ($n\gg0$)。
\item 三連続の一次Homと二次合成核から作る空間$D_3(\mathcal B_0)$は二十次元。
\end{enumerate}
このとき$\mathcal B\not\simeq\operatorname{coh}\mathcal A$である。
\end{proposition}
\begin{proof}
主論文の非表示定理の証明は、この四条件と
$\operatorname{coh}\mathcal A$側の構造だけを使う。
同値による$\mathcal B_0$の像は局所射影division rank一exceptionalとなり、
分類と$D_3$の次元比較が十九と二十の矛盾を与える。
\end{proof}

任意の非周期的な標準AS正則$\mathbb Z$-algebraとその正規三次商が、
この四条件を全て満たすという主張は、現稿では証明していない。
特に必要なのは商のdualizing bimoduleとtautological対象列へのSerre作用である。
Mori--Nyman\cite{MoriNyman}のv3にはcorrigendumがあり、
旧Theorem~7.8は追加のNakayama bimodule条件なしには使えない。
Corollary~7.6のrowwise局所コホモロジー公式だけから
全bimodule双対性を結論しない。

次の具体的方針は、両側14641型Koszul dualのFrobenius構造から
周期四のNakayama bimoduleを取り出し、
正規三次イデアルのshift三とnormal-divisor adjunctionを組み合わせることである。
これは今後証明すべき技術的補題であり、現在のgraded定理の依存結果ではない。

\section{使える補助事実と避けるべき推論}
\subsection{Hochschild homology}
次数三の例は
\[
(\dim HH_{-1},\dim HH_0,\dim HH_1)=(2,4,2)
\]
を持ち、full exceptional collectionを持たない。
これはKuznetsov\cite[Theorem~8.12]{KuznetsovHH}と、
種数三の曲線の非自明2-torsion線束のcohomologyから計算できる。
ただし、非零の$HH_{\pm1}$それ自体を、
全てのAS四次元三次超曲面に対する障害としては用いない。
本稿の排除は二次包絡と候補構造対象の分類に基づく。

\subsection{数値的群の計算}
conic bundle $p:X\to Z$は滑らかな三次元多様体である。
連結判別二重被覆により$p^{-1}(D)$は既約であり、
一般conicへのPicard群制限から$\rho(X)=3$となる。
$H^{>0}(X,\mathcal O_X)=0$とLefschetzの定理から
$\rk K_{\mathrm{num}}(X)=8$、半直交分解から
$\rk K_{\mathrm{num}}(\mathcal A)=4$を得る。
$X$をambientのample divisorと仮定する必要はない。

基底$(\mathcal A,\mathcal P,\mathcal A(C),\mathcal A(f))$で
\[
G=\begin{pmatrix}
1&4&0&2\\1&1&-2&0\\-2&-1&1&-5\\0&1&-5&1
\end{pmatrix},\qquad\det G=100.
\]
局所射影rank一$F$について
$\det F=\det\mathcal A\otimes N^2$、$N=aC+bf$とすると、
exceptionalityは
\[
\chi(\mathcal A,F)=1+\frac{N^2+NL}{2}
\]
を与える。整数性から$b$は偶数である。
従って$N$のmodulo $2\operatorname{Pic}Z$の類は$0$または$L$になり、
$\mathcal A(M)$、$\mathcal P(M)$が全候補となる。
これは単なる数値的分類に止まらず、行列式とSerre双対性で実際の同型に上がる。

\section{一次資料と既知の範囲}
\begin{itemize}
\item Chan--Ingalls\cite[Sections~2--3]{ChanIngalls2012}：
線束値二次形式、Clifford偶部と奇部、conic bundleとの関係。
Section~8の具体的minimal del Pezzo例は主として中心$\mathbb P^2$上を扱う。
\item Nasr\cite[arXiv v1, Definition~2.1, Theorems~1.1(2), 2.16]{Nasr2020}：
本例の$D=4H-C$、分岐指数二は既知の一点blow up型に含まれる。
orderの有理標準因子は$K_Z+\tfrac12D$である。
\item Kuznetsov\cite[Theorem~4.2]{KuznetsovQuadric}：
flat quadric fibrationの半直交分解。
本例では基底と全空間の滑らかさを先に確認して適用する。
\item Yekutieli--Zhang\cite[Theorem~4.2, Corollary~4.3]{YZ1997}：
左右noetherianなconnected AS-Gorenstein代数のbalanced dualizing complex、
局所コホモロジーの右有界性、Serre双対性。
原著のpp.\ 705--707を確認した。
\item Artin--Zhang\cite[Theorem~4.5]{AZ1994}：
ample pairの再構成に用いる。
書誌・出版社ページを確認した。今回、当該原著の全文再取得はできていない。
使用する再構成の標準的な定式化と、本例における仮定は本文で確認した。
\item Cassidy--Vancliff\cite[Theorem~1.1]{CV2010}：
basepoint-free な三つの二次形式から作る graded Clifford algebra の AS 正則性。
次数二の別 ambient の構成で使う。著者公開 PDF の該当定理を確認した。
\item Chirvasitu--Kanda--Smith\cite[arXiv v3, Propositions~1.6, 2.1]{CKS2019}：
連結次数付き twisted Calabi--Yau と AS 正則性の関係、および正則正規拡大。
DOI は \texttt{10.1090/tran/7672} である。
\item Chen--Kirkman--Zhang\cite[Lemma~1.1]{CKZ2017} と
Liu--Wang--Wu\cite[Theorem~3.3]{LWW2014}：
down-up algebra と Ore 拡大の正則性。
次数一の最小生成次数 $(1,1,3)$ の ambient に適用する。
\item Bernardara--Macr\`i--Mehrotra--Stellari
\cite[Propositions~2.9(ii), 2.12]{BMMS2012}：
cubic-threefold Clifford 圏の数値的 Grothendieck 群の整数基底。
次数一の対象分類に用いる。原著 PDF の両命題を確認した。
\end{itemize}

Chan--Ingalls の $F_{5-}^{2}$ 型は arXiv v1 では Section~9、
最終著者版では Section~8 にある。本研究の次数一の分解と一致する。
Kuznetsov の quadric 論文は出版版
\emph{Adv. Math.} 218 (2008), 1340--1369 を引用し、
Theorem~4.2 が p.~1357 にあることを確認した。
Hochschild 論文は確認できた arXiv:0904.4330v1 を用い、未確認の雑誌情報を補わない。
Nasr の書誌情報は出版版によるが、定義・定理番号は確認した arXiv v1 を明記する。
Serre vanishing は Yekutieli--Zhang の局所双対性から本文で直接導き、
今回原著本文を取得できなかった AZ Corollary~7.5 には依存させない。

比較の出発点となる $\mathcal H_3$ の定義は、2026年11月19日に次の原稿で確認した。
その二次包絡の $(1,4,6,4,1)$ 条件と本研究の十九次元の計算は直接比較できる。
本例では、元の weak な内在的条件に加えて、smooth proper 性と
全整数 $n$ に対する $\operatorname{Ext}^1(\mathcal O,T^n\mathcal O)=0$ も証明している。
\begin{quote}\small
\url{https://github.com/uedakazushi/nc-cubic-moduli}\par
\texttt{papers/weak\_cubic\_axioms.tex}\par
blob SHA: \texttt{7190004e0914225d186cb6dd383a5ed8d85af540}
\end{quote}
主論文には必要な定義を全て再掲してあり、再開に外部原稿の取得は不要である。

\section{再現資産と検証状態}
\path{scripts/verify_all.py} は、元の有限群sector計算、Clifford次数一・二の計算、
次数三の滑らかさ・二次関係の計算と六標識の整数行空間の計算をまとめて実行する。
従来のCliffordの五計算は Python 3.12.14、SymPy 1.14.0 での成功記録があり、
現行版の一括再実行とは区別する。
六標識の新しいコードはPython標準ライブラリで再実行し成功を確認した。
計算では有理数・整数演算を用い、浮動小数点近似を用いない。
行列証明書は \path{certificates/degree3_quadratic.json} として保持する。
通常の検証はこれを読み取り、変更しない。
四十五次小行列式が $512$、三本の整数核ベクトルが独立かつ実際に核に属することを確認する。
証明書を再生成する場合だけ
\texttt{python3 scripts/verify\_degree3\_relations.py --regenerate} を実行する。

\path{scripts/build.py} は LuaLaTeX、BibTeX、latexmk で二つの文書を組版する。
必要な依存関係と VS Code / Cursor の操作は README に記す。
\path{scripts/make_bundle.py} は現在のソースと主論文と研究ノートのPDF だけをまとめ、
ログ、キャッシュ、旧版、入れ子 ZIP を含めない。
実行ログを保管する代わりに入力・コード・必要な数学的証明書を保持する。


\section{量子Fermat点と非CY符号型族の現行の要点}
本節は統合先に既にあった二つのH3例のうち、
主論文に十分な証明がある事項の参照先と、再検証時の注意点を整理する。
量子Fermatは三乗根の量子多項式環の中心三次商で、
離散torsion付き$(\mu_3)^2$商によるMorita表示と
$HH_{-1,0,1}=(2,5,2)$ を与える。
非CY符号型では$q_{01}=1$、他の五つの交換係数は$-1$、
$G=f_3(x_0,x_1)+x_0x_2^2+x_1x_3^2$、滑らかさ条件は
$a d\Delta(f_3)\ne0$である。
この族は固定ambientで一般に二次元のpointed同型類を持ち、
$HH_{-1,0,1}=(2,4,2)$である。

符号型では$\Gamma=(\mathbb Z/2)^3$、
$\deg(x_0,x_1,x_2,x_3)=(e_0,e_0,e_1,e_2)$、
cocycle $\nu(u,v)=(-1)^{u_1v_0+u_2v_0+u_2v_1}$を使う。
交換形式のradicalは$e_0+e_1+e_2$なので、
twisted group algebraは$M_2(\mathbb C)\oplus M_2(\mathbb C)$であり、
$M_4(\mathbb C)$と取り違えない。
その二成分を交換する$H=\Gamma^\vee$の安定化群は
$H_0=\{(a,b,c):a+b+c=0\}$である。
非自明なsector $g_A=(0,1,1)$の固定集合は
$\mathbb P^1$と三点、$g_B=(1,0,1)$と$g_C=(1,1,0)$は
それぞれ楕円曲線と一点を含む。
$H^2$のtraceは$(7,3,-1,-1)$、$H^2$不変部分のrankは二である。
twisted HKRの非零奇数次部分は二つの楕円sectorに由来する。
詳細な導来Morita還元とsector算定の証明は主論文の
量子Fermat・符号型節を参照し、計算確認は
\path{computations/verify_all.py}に置く。

本研究の二つの非存在概念は異なる。
量子Fermat・符号型は標準四次元AS regular ambientの
三次超曲面として表示されるがOU型AS surfaceでない。
次数三のClifford例は、反標準環の二次包絡が
標準四次元AS regular環でなく、裸のAbel圏としても
その種の三次超曲面の$\operatorname{qgr}$と同値にならない。
次数一・二Clifford例では指定最小生成次数の非表示と
別ambientの存在を両立させる。

\bibliographystyle{amsplain}
\bibliography{references}
\end{document}
